\documentclass[spanish,11pt,a4paper]{article}
\usepackage[margin=1in]{geometry}
\usepackage{fontspec}
\setmainfont{lmroman10-regular.otf}[BoldFont=lmroman10-bold.otf,ItalicFont=lmroman10-italic.otf,BoldItalicFont=lmroman10-bolditalic.otf]
\setsansfont{lmsans10-regular.otf}[BoldFont=lmsans10-bold.otf]
\setmonofont{lmmono10-regular.otf}[Scale=MatchLowercase]
\newfontfamily\LeanSymbols{seguisym.ttf}[Path=typeset_runtime/]
\newlength{\LeanCharWidth}
\usepackage{amsmath,amssymb}
\usepackage[spanish,es-noshorthands,es-tabla]{babel}
\usepackage{microtype,parskip}
\usepackage{graphicx,float,caption}
\usepackage{longtable,booktabs,array,calc}
\usepackage{xcolor}
\usepackage{xurl}
\usepackage{hyperref,bookmark}
\usepackage{etoolbox,needspace}
\setcounter{secnumdepth}{-\maxdimen}
\setlength{\emergencystretch}{3em}
\providecommand{\tightlist}{\setlength{\itemsep}{0pt}\setlength{\parskip}{0pt}}
\newcounter{none}
\AtBeginEnvironment{longtable}{\Needspace{8\baselineskip}\small}
\pretocmd{\section}{\Needspace{5\baselineskip}}{}{}
\pretocmd{\subsection}{\Needspace{4\baselineskip}}{}{}
\pretocmd{\subsubsection}{\Needspace{4\baselineskip}}{}{}
\captionsetup{font=small,labelfont=bf}
\urlstyle{same}
\providecommand{\pandocbounded}[1]{#1}
\hypersetup{pdftitle={Paper III: formalización y reproducibilidad},pdfauthor={Juan Pablo Traverso Gianini},pdflang={es-ES},colorlinks=true,linkcolor=blue,urlcolor=blue,citecolor=blue}
\title{Paper III: formalización y reproducibilidad}
\author{Juan Pablo Traverso Gianini\\\small Investigador independiente, Santiago, Chile\\\small \href{mailto:jtraverso@gmail.com}{jtraverso@gmail.com}\\\small ORCID: 0009-0003-6068-4096}
\date{}
\begin{document}
\raggedbottom
\maketitle
\textbf{Paper III · Formal companion v1.6}\\
19 de octubre de 2026

Este documento acompaña al artículo matemático. Conserva las
declaraciones formales, las precisiones de alcance y la evidencia
computacional que antes figuraban en el artículo. El congelado de
fuentes y su verificación registrada no cambian. No se ejecutó Lean/Lake
para esta edición. Los números históricos de secciones y tablas que
siguen remiten al artículo anterior; se conservan como localizadores y
no se renumeran como resultados nuevos de verificación. Las referencias
originales {[}1{]}, {[}2{]}, etc. figuran al final.

\href{https://github.com/jtraverso/erdos-81-chordal-clique-partitions/tree/eb335845bfa23375d5040ea31fc4f4b366346287/preprints/PAPER_III/05_formalization}{Frozen
sources and historical evidence / Fuentes congeladas y evidencias
históricas} ·
\href{https://github.com/jtraverso/erdos-81-chordal-clique-partitions/tree/eb335845bfa23375d5040ea31fc4f4b366346287/preprints/PAPER_III/01_manuscript}{Previous
article / Artículo anterior}

\subsection{11.6 Perímetro de verificación formal y estado
actual}\label{peruxedmetro-de-verificaciuxf3n-formal-y-estado-actual}

La discusión de la formalización se concentra en esta sección y en el
registro de release de la Sección 13. El snapshot formal actual separa
el ensamblaje final incondicional de sus interfaces downstream. La raíz
pública agregada \nolinkurl{PaperIII} importa tanto
\nolinkurl{PaperIII.Theorem_1_1_Final} como \nolinkurl{PaperIII.PublicAPI}.
El módulo final descarga las dos interfaces usadas por el artículo y
exporta el Teorema 1.1 y el Corolario 1.2.
\nolinkurl{PaperIII.CanonicalTrianglePacking} fija \nolinkurl{PaperIII.nu3} y
\nolinkurl{PaperIII.nu3Star} como cantidades orientadas al manuscrito y
demuestra su igualdad con las implementaciones mediante hipergrafos de
aristas utilizadas por AX1; también demuestra
\nolinkurl{PaperIII.tau3Star\ =\ PaperIII.nu3Star} e identifica la interfaz
AX1 del lado de coberturas con el enunciado del lado de empaquetamientos
de la Sección 2. \nolinkurl{PaperIII.PaperImprovements} contiene los
certificados orientados al artículo; \nolinkurl{PaperIII.PublicAPI},
\nolinkurl{PaperIII.OfPartition} y \nolinkurl{PaperIII.Obstructions} exponen
las interfaces downstream. El gate canónico dirigido informa la
siguiente huella de axiomas para estas superficies:

\begin{quote}
{\LeanSymbols\fontsize{8}{10}\selectfont
\noindent \nolinkurl{[propext,}\hspace{0.45em}\nolinkurl{Classical.choice,}\hspace{0.45em}\nolinkurl{Quot.sound]}\par
}
\end{quote}

El ensamblaje final usa puentes explícitos hacia las formalizaciones de
Haxell--Rödl/Yuster y de descomposición densa. Los teoremas asintóticos
citados se mantienen en el manuscrito por su procedencia matemática. El
build registrado del lado del autor para la raíz pública completó 8.455
jobs con exit code 0; el build posterior de las raíces de consulta
completó 8.444 jobs con exit code 0. Ese proceso del autor se reanudó
después del reinicio de una aplicación. Por separado, la auditoría
externa realizó una reproducción ininterrumpida desde la ausencia del
directorio de build del proyecto: la raíz pública completó 8.455 jobs,
las raíces de consulta completaron 8.444 jobs y los ocho archivos de
consulta de axiomas terminaron correctamente. El archivo formal
inmutable está sellado, y los gates Lean internos y externos pasan
dentro del alcance declarado en la Sección 13.

El Apéndice E ofrece una selección curada de certificados reutilizables
y de nivel paper. El inventario completo al nivel de implementación
permanece en la fuente formal, en lugar de duplicarse en el manuscrito.

\textbf{Tabla 2. Superficies de formalización.} Los estados de esta
tabla son registros del freeze local sellado; no sustituyen la
reproducción independiente.

{\def\LTcaptype{none} % do not increment counter
\begin{longtable}[]{@{}
  >{\raggedright\arraybackslash}p{(\linewidth - 4\tabcolsep) * \real{0.3333}}
  >{\raggedright\arraybackslash}p{(\linewidth - 4\tabcolsep) * \real{0.3333}}
  >{\raggedright\arraybackslash}p{(\linewidth - 4\tabcolsep) * \real{0.3333}}@{}}
\toprule\noalign{}
\begin{minipage}[b]{\linewidth}\raggedright
Superficie
\end{minipage} & \begin{minipage}[b]{\linewidth}\raggedright
Propósito
\end{minipage} & \begin{minipage}[b]{\linewidth}\raggedright
Estado local actual
\end{minipage} \\
\midrule\noalign{}
\endhead
\bottomrule\noalign{}
\endlastfoot
\nolinkurl{PaperIII} & raíz agregada/de compatibilidad & el build de la
raíz pública completó 8.455 jobs; exit 0; importa el teorema final y la
API pública \\
\nolinkurl{PaperIII.MainNibble} & ensamblaje del teorema y los corolarios
mediante interfaces AX1/AX2 explícitas & compilada en la clausura
registrada \\
\nolinkurl{PaperIII.Theorem_1_1_Final} & descarga incondicional de
AX1/AX2 y teorema final & compilada; la huella del teorema final es
exclusivamente fundacional \\
\nolinkurl{PaperIII.CanonicalTrianglePacking} & puentes canónicos para
\nolinkurl{nu3}, \nolinkurl{nu3Star}, \nolinkurl{tau3Star} y el enunciado AX1 &
compilada; se registraron los puentes exactos del target \\
\nolinkurl{PaperIII.CanonicalTrianglePackingGate} & gate de enunciados y
axiomas para las interfaces canónicas & exit 0; huellas exclusivamente
fundacionales registradas \\
\nolinkurl{PaperIII.PaperImprovements} y
\nolinkurl{PaperIII.PaperImprovementsGate} & corolarios orientados al paper
y certificados auxiliares seleccionados & compilados en la corrida
consolidada de gates \\
\nolinkurl{PaperIII.PublicAPI}, \nolinkurl{PaperIII.OfPartition},
\nolinkurl{PaperIII.Obstructions} & API downstream, transporte de partición
split y certificados de obstrucción & compilados en la clausura
registrada \\
\nolinkurl{FREEZE_REPORT.md} y \nolinkurl{ESCAPE_HATCH_ASSESSMENT.md} &
evaluación inmutable del build, los axiomas y los tokens crudos &
generados e incluidos en el archivo formal sellado \\
\end{longtable}
}

Al nivel de \nolinkurl{SimpleGraph},
\nolinkurl{Theorem_1_1_of_splitPartition_uncond} acepta un grafo
finito provisto de una partición clique/independiente. La construcción
\nolinkurl{SplitGraph.ofPartition} se conecta con el grafo original
mediante \nolinkurl{ofPartitionIso}; el número de vértices, el número de
aristas, \(\nu_3\) y \(\Phi\) se transportan a través de la invariancia
por isomorfismos de grafos, incluida \nolinkurl{nu3_congr_of_iso}. Otras
formas del teorema especializadas por tipos permanecen en
\nolinkurl{PaperIII.PublicAPI} para uso downstream, pero no son resultados
matemáticos del artículo y no se catalogan aquí.

La compilación dirigida del módulo canónico, los builds de las raíces
pública y de consulta, los ocho archivos compatibles de consulta de
axiomas, la evaluación de escape hatches, los manifiestos inmutables de
fuentes y la auditoría interna pasan sus controles del lado del autor.
La auditoría externa reprodujo independientemente las raíces pública y
de consulta y verificó las superficies de axiomas designadas contra los
mismos bytes congelados.

Ningún cálculo ni script de regresión es premisa de la demostración. Las
comprobaciones de aritmética exacta y programación entera siguen siendo
auditorías suplementarias. La revisión adversarial completada y la
evaluación de arte previo acotada por su corpus no sustituyen una
revisión humana por pares ni una determinación especializada de
prioridad; el etiquetado público es una operación de release al nivel
del repositorio.

\subsection{11.7 Componentes reutilizables de la
formalización}\label{componentes-reutilizables-de-la-formalizaciuxf3n}

La formalización también produjo una biblioteca independiente
\nolinkurl{Contrib.*}, cuyos módulos importan únicamente Mathlib. Cuatro
grupos resultan especialmente pertinentes más allá de este artículo:

\begin{itemize}
\tightlist
\item
  dualidad fuerte finita para LP de empaquetamiento/cobertura mediante
  un argumento de Farkas cónico (\nolinkurl{Contrib.LPStrongDuality});
\item
  desigualdades de concentración sensibles a la varianza y colas de
  martingalas (\nolinkurl{Contrib.BennettBernstein},
  \nolinkurl{Contrib.HoeffdingUpper}, \nolinkurl{Contrib.ConditionalHoeffding}
  y \nolinkurl{Contrib.FreedmanBernstein});
\item
  resultados de tipo Dirac para hamiltonicidad y matchings
  (\nolinkurl{Contrib.DiracHamiltonian} y
  \nolinkurl{Contrib.SpreadMatchingDirac});
\item
  componentes finitos de flujo máximo/corte mínimo y de teoría de
  diseños (\nolinkurl{Contrib.MaxFlowMinCut} y \nolinkurl{Contrib.FanoSTS}).
\end{itemize}

Estos módulos son subproductos reutilizables de la formalización, no
premisas del Teorema 1.1. Se reservan para una revisión separada y una
posible contribución upstream. El presente manuscrito no afirma que sus
APIs o sus enunciados individuales sean nuevos respecto de la biblioteca
Mathlib vigente; esa comparación exige una revisión separada de
literatura y de la API orientada a upstream. El Apéndice E.6 registra
solo los componentes más cercanos a la demostración matemática.

\begin{center}\rule{0.5\linewidth}{0.5pt}\end{center}

\subsection{13. Reproducibilidad}\label{reproducibilidad}

La fuente matemática es este manuscrito. Las superficies formales
orientadas al manuscrito son \nolinkurl{PaperIII.Theorem_1_1_Final},
\nolinkurl{PaperIII.CanonicalTrianglePacking}, \nolinkurl{PaperIII.PublicAPI},
\nolinkurl{PaperIII.OfPartition} y \nolinkurl{PaperIII.Obstructions}.
\nolinkurl{PaperIII.PaperImprovementsGate} es el gate de subproductos, y
\nolinkurl{PaperIII.CanonicalTrianglePackingGate} comprueba la
correspondencia exacta entre los parámetros canónicos de empaquetamiento
y el modelo Nibble/Yuster usado para descargar AX1.

El desarrollo formal está sellado como
\nolinkurl{PAPER_III_lean_v1.4_freeze.zip}, con SHA-256
\nolinkurl{79ee24c38fd776bc2585a0c3c996e30817f0829fc5064463bdbde0fa2d3d7104}.
Contiene 704 archivos fuente Lean y los tres archivos fijados de
Lake/toolchain enumerados en el manifiesto de 707 entradas. La
verificación registrada usa Lean 4.28.0 y Mathlib v4.28.0 en el commit
\nolinkurl{8f9d9cff6bd728b17a24e163c9402775d9e6a365}.

El registro local de verificación disponible se resume a continuación.

\textbf{Tabla 3. Registro local sellado de verificación.}

{\def\LTcaptype{none} % do not increment counter
\begin{longtable}[]{@{}
  >{\raggedright\arraybackslash}p{(\linewidth - 4\tabcolsep) * \real{0.3333}}
  >{\raggedright\arraybackslash}p{(\linewidth - 4\tabcolsep) * \real{0.3333}}
  >{\raggedright\arraybackslash}p{(\linewidth - 4\tabcolsep) * \real{0.3333}}@{}}
\toprule\noalign{}
\begin{minipage}[b]{\linewidth}\raggedright
Comprobación
\end{minipage} & \begin{minipage}[b]{\linewidth}\raggedright
Log o comando
\end{minipage} & \begin{minipage}[b]{\linewidth}\raggedright
Resultado actual
\end{minipage} \\
\midrule\noalign{}
\endhead
\bottomrule\noalign{}
\endlastfoot
build de la raíz pública agregada &
\nolinkurl{gate_logs/run_2026-08-22_CCS_NOTEBOOK456_resumed/logs/02_build_public_root_clean.log}
& exit 0; 8.455 jobs completados; proceso de origen limpio reanudado
después del reinicio de una aplicación \\
completación de las raíces de consulta &
\nolinkurl{gate_logs/run_2026-08-22_CCS_NOTEBOOK456_resumed/logs/03_build_query_roots_incremental.log}
& exit 0; 8.444 jobs completados \\
gate canónico de enunciados y axiomas &
\nolinkurl{PaperIII.CanonicalTrianglePackingGate} & exit 0; comprobados los
puentes integrales y fraccionales del target \\
registro completo de axiomas por teorema & ocho archivos
\nolinkurl{FreezeAxioms*.lean} & exit 0; solo \nolinkurl{propext},
\nolinkurl{Classical.choice} y \nolinkurl{Quot.sound} \\
regresión cuantitativa del cierre Lean &
\nolinkurl{FreezeAxiomsAuditClosure.lean} y
\nolinkurl{FreezeAxiomsAX1Closure.lean} & todas las superficies nombradas
están presentes con huellas exclusivamente fundacionales \\
evaluación de escape hatches & escaneo de tokens crudos y controles
enfocados de build/huellas & sin warning activo de \nolinkurl{sorry}, sin
\nolinkurl{native_decide} activo y sin axioma de proyecto en una
superficie canónica \\
sello inmutable & manifiestos de fuentes/paquete, ZIP y sidecar SHA-256
& 707 entradas fuente; hash del archivo verificado \\
\end{longtable}
}

La verificación registrada usó un único snapshot de fuentes sin cambios.
Su secuencia de build fue

\begin{quote}
{\LeanSymbols\fontsize{8}{10}\selectfont
\noindent \nolinkurl{lake}\hspace{0.45em}\nolinkurl{build}\hspace{0.45em}\nolinkurl{PaperIII}\par
\noindent \nolinkurl{lake}\hspace{0.45em}\nolinkurl{build}\hspace{0.45em}\nolinkurl{BKLO.MainDenseUnconditional}\hspace{0.45em}\nolinkurl{Nibble.AX1Closed}\hspace{0.45em}\nolinkurl{PaperIII.CanonicalTrianglePacking}\hspace{0.45em}\nolinkurl{PaperIII.Obstructions}\hspace{0.45em}\nolinkurl{PaperIII.PaperImprovementsGate}\hspace{0.45em}\nolinkurl{PaperIII.PublicAPI}\hspace{0.45em}\nolinkurl{PaperIII.Theorem_1_1_Final}\par
\noindent \nolinkurl{#print}\hspace{0.45em}\nolinkurl{axioms}\hspace{0.45em}\nolinkurl{PaperIII.Theorem_1_1}\par
\noindent \nolinkurl{#print}\hspace{0.45em}\nolinkurl{axioms}\hspace{0.45em}\nolinkurl{PaperIII.AX1_holds}\par
\noindent \nolinkurl{#print}\hspace{0.45em}\nolinkurl{axioms}\hspace{0.45em}\nolinkurl{PaperIII.AX2_holds}\par
}
\end{quote}

El proyecto del lado del autor comenzó sin directorio \nolinkurl{.lake} ni
objetos compilados del proyecto. La aplicación de escritorio se reinició
durante el primer proceso de la raíz pública; el proyecto, sin cambios,
se reanudó entonces de manera incremental y terminó correctamente. Por
ello, el freeze inmutable registra
\nolinkurl{PASS_CLEAN_ORIGIN_RESUMED}, no una única invocación
ininterrumpida del autor. Conserva esa declaración, la salida verbatim,
los identificadores de componentes, el manifiesto de fuentes, el
manifiesto del archivo y los inventarios SHA-256. La auditoría externa
aportó después una reproducción ininterrumpida desde la ausencia del
directorio de build del proyecto, con las revisiones fijadas de las
dependencias: 8.455 jobs de la raíz pública y 8.444 jobs de las raíces
de consulta terminaron con exit code 0, y pasaron los ocho archivos de
consulta de axiomas. Esos controles al nivel de declaraciones establecen
conformidad del build y huellas de axiomas; no reemplazan una
rederivación matemática independiente.

Dos módulos legacy archivados, \nolinkurl{Ax2/PartB/Axioms.lean} y
\nolinkurl{Ax2/PartA/Wlog.lean}, contienen declaraciones de axiomas locales
del proyecto. Ningún archivo fuente importa esos módulos y ninguna de
esas declaraciones aparece en la huella de dependencias de una
superficie canónica de release. Por tanto, el claim formal preciso es al
nivel de declaraciones: el teorema canónico y las interfaces de release
no tienen axiomas matemáticos locales del proyecto; no se afirma que el
árbol histórico completo carezca del token \nolinkurl{axiom}.

La demostración misma permanece en el manuscrito. Los scripts de
regresión y las auditorías computacionales no son hipótesis, reducciones
ni pasos de la demostración.

El paquete del repositorio canónico contiene la fuente Lean completa, el
toolchain, el manifiesto, los logs de build y axiomas, las fuentes del
manuscrito y los inventarios SHA-256. El tag y el commit públicos
vinculan esos archivos al nivel del repositorio; el manuscrito no afirma
dentro de su texto un identificador de commit.

Los materiales escritos se destinan a distribución bajo CC BY-NC 4.0. El
código de verificación y regresión propiedad del autor se destina a
distribución bajo PolyForm Noncommercial License 1.0.0; las dependencias
de terceros conservan sus propias licencias.

\textbf{Límite de la revisión:} la reproducción independiente y la
auditoría adversarial externa están completas, y la revisión de novedad
está acotada por su corpus. No se han realizado una revisión humana por
pares ni una determinación especializada de prioridad. El tag y el
commit públicos son metadatos del repositorio y no premisas del teorema.

\begin{center}\rule{0.5\linewidth}{0.5pt}\end{center}

\subsection{Apéndice C. Auditorías
computacionales}\label{apuxe9ndice-c.-auditoruxedas-computacionales}

\subsection{C.1 LP exacto de perfil
común}\label{c.1-lp-exacto-de-perfil-comuxfan}

\nolinkurl{verify_common_profile_lp.py} enumera los vértices del
poliedro simetrizado de coberturas mediante eliminación gaussiana
racional exacta y compara el óptimo con (3.5).

\subsection{C.2 Margen fraccional
exacto}\label{c.2-margen-fraccional-exacto}

\nolinkurl{verify_fractional_margin.py} comprueba Teorema 4.2 para

\[
3\le p\le80,
\qquad
1\le q\le2p,
\qquad
0\le d\le p
\]

usando aritmética racional exacta.

\subsection{C.3 Auditorías ILP
pequeñas}\label{c.3-auditoruxedas-ilp-pequeuxf1as}

\nolinkurl{verify_factor_rounding.py} y
\nolinkurl{verify_shifted_center.py} calculan empaquetamientos integrales
exactos de triángulos en instancias pequeñas mediante un ILP binario con
capacidades de aristas. Verifican las cotas superiores enunciadas para
\(\Phi(G)\).

\subsection{C.4 Polarización}\label{c.4-polarizaciuxf3n}

\nolinkurl{verify_polarization.py} ejecuta comprobaciones exhaustivas de
perfiles pequeños y comprobaciones aleatorias con enteros exactos de
Lema 6.1.

\subsection{C.5 Divisibilidad}\label{c.5-divisibilidad}

\nolinkurl{verify_divisibility_correction.py} verifica exhaustivamente la
construcción de paridad sobre caminos hasta orden dieciocho, comprueba
exactamente el umbral fijo de grado mínimo y sigue la corrección
completa de paridad y módulo tres sobre residuales densos aleatorios.

\begin{center}\rule{0.5\linewidth}{0.5pt}\end{center}

\subsubsection{\texorpdfstring{Observación 2.3\(^{\prime}\) ---
Necesidad de una hipótesis de
densidad}{Observación 2.3\^{}\{\textbackslash prime\} --- Necesidad de una hipótesis de densidad}}\label{observaciuxf3n-2.3prime-necesidad-de-una-hipuxf3tesis-de-densidad}

\textbf{Certificados Lean.}
\nolinkurl{PaperIII.ax2_divisibility_degree_insufficient},
\nolinkurl{PaperIII.ax2_density_necessary_K7_minus_two_triangles},
reexportados mediante \nolinkurl{PaperIII.Obstructions}.

\subsubsection{Corolario 10.2b --- Término cuadrático afilado y
coeficiente lineal
forzado}\label{corolario-10.2b-tuxe9rmino-cuadruxe1tico-afilado-y-coeficiente-lineal-forzado}

\textbf{Certificado Lean.} \nolinkurl{PaperIII.Corollary_1_2_sharp},
ensamblado a partir de \nolinkurl{PaperIII.Corollary_1_2},
\nolinkurl{PaperIII.Byproduct_leading_constant_forced} y
\nolinkurl{PaperIII.Byproduct_completeSplit_cp_sharp}. La declaración
\nolinkurl{Byproduct_leading_constant_forced} certifica la afirmación
sobre el coeficiente lineal; la sharpness cuadrática se sigue del valor
exacto de la familia.

\subsubsection{Corolario E.1.1 --- Empaquetamiento por asignación de
factores}\label{corolario-e.1.1-empaquetamiento-por-asignaciuxf3n-de-factores}

\textbf{Certificado Lean.}
\nolinkurl{Byproduct_factorization_assignment_packing}.

\subsubsection{Corolario E.1.2 --- Cota inferior por doble
factorización}\label{corolario-e.1.2-cota-inferior-por-doble-factorizaciuxf3n}

\textbf{Certificado Lean.}
\nolinkurl{Byproduct_double_factorization_packing}.

\subsubsection{Corolario E.1.3 --- Cota de empaquetamiento con ganancia
reservada}\label{corolario-e.1.3-cota-de-empaquetamiento-con-ganancia-reservada}

\textbf{Certificado Lean.}
\nolinkurl{Byproduct_reserved_gain_packing_bound_subset}.

\subsubsection{\texorpdfstring{Corolario E.2.1 --- Constante \(3/2\) del
corredor
corto}{Corolario E.2.1 --- Constante 3/2 del corredor corto}}\label{corolario-e.2.1-constante-32-del-corredor-corto}

\textbf{Certificado Lean.}
\nolinkurl{Byproduct_short_corridor_constant_three_halves}.

\subsubsection{Corolario E.2.2 --- Puente de empaquetamientos a
particiones}\label{corolario-e.2.2-puente-de-empaquetamientos-a-particiones}

\textbf{Certificados Lean.} \nolinkurl{Byproduct_cp_le_Phi},
\nolinkurl{Byproduct_low_corridor_cp_three_halves},
\nolinkurl{Byproduct_effective_low_corridor_cp_bound}.

\subsubsection{Corolario E.3.1 --- Óptimo fraccional
exacto}\label{corolario-e.3.1-uxf3ptimo-fraccional-exacto}

\textbf{Certificado Lean.}
\nolinkurl{Byproduct_commonProfile_tau3Star_eq_F}.

\subsubsection{Corolario E.3.2 --- Regiones de
fase}\label{corolario-e.3.2-regiones-de-fase}

\textbf{Certificados Lean.}
\nolinkurl{Byproduct_commonProfile_tau3Star_eq_uniform_region},
\nolinkurl{Byproduct_commonProfile_tau3Star_eq_separated_region},
\nolinkurl{Byproduct_commonProfile_tau3Star_eq_hot_region}.

\subsubsection{Corolario E.4.1 --- Cota inferior para la holgura
fraccional}\label{corolario-e.4.1-cota-inferior-para-la-holgura-fraccional}

\textbf{Certificado Lean.}
\nolinkurl{Byproduct_unified_fractional_margin}.

\subsubsection{Corolario E.5.1 --- Valor exacto y forma
normal}\label{corolario-e.5.1-valor-exacto-y-forma-normal}

\textbf{Certificados Lean.}
\nolinkurl{Byproduct_completeSplit_cp_exact_value},
\nolinkurl{Byproduct_completeSplit_cp_sharp}.

\subsubsection{Corolario E.5.2 --- Coeficiente lineal
forzado}\label{corolario-e.5.2-coeficiente-lineal-forzado}

\textbf{Certificados Lean.}
\nolinkurl{Byproduct_leading_constant_forced}, empaquetado con la cota
superior y la familia exacta como \nolinkurl{Corollary_1_2_sharp}.

\subsubsection{E.6.1 Dualidad LP finita
empaquetamiento/cobertura}\label{e.6.1-dualidad-lp-finita-empaquetamientocobertura}

\textbf{Certificados Lean.} \nolinkurl{Byproduct_triangle_LP_duality};
forma independiente \nolinkurl{Contrib.LPDuality.lp_strong_duality} en
\nolinkurl{Contrib.LPStrongDuality}.

\subsubsection{E.6.2 Matchings de tipo Dirac y corrección
densa}\label{e.6.2-matchings-de-tipo-dirac-y-correcciuxf3n-densa}

\textbf{Certificados Lean.}
\nolinkurl{Byproduct_dirac_perfect_matching_even},
\nolinkurl{Byproduct_dirac_near_perfect_matching},
\nolinkurl{Byproduct_dense_divisible_correction_edges}; forma
independiente de matching en \nolinkurl{Contrib.SpreadMatchingDirac}.

\subsubsection{E.6.3 Interfaz uniforme de
nibble}\label{e.6.3-interfaz-uniforme-de-nibble}

\textbf{Certificado Lean.}
\nolinkurl{Byproduct_almostPerfectMatching_uniform}.

\subsubsection{Corolario E.7.1 --- Realización y cota fraccional por
aristas}\label{corolario-e.7.1-realizaciuxf3n-y-cota-fraccional-por-aristas}

\textbf{Certificados Lean.} \nolinkurl{Byproduct_nu3_achieved},
\nolinkurl{Byproduct_nu3star_le_edges_div_three}.

\subsubsection{Corolario E.7.2 --- Forma de brecha de aristas del
majority
nibble}\label{corolario-e.7.2-forma-de-brecha-de-aristas-del-majority-nibble}

\textbf{Certificado Lean.}
\nolinkurl{Byproduct_nu3star_sub_nu3_le_edge_gap_most}.

\subsubsection{Apéndice E. Corolarios seleccionados derivados de la
formalización}\label{apuxe9ndice-e.-corolarios-seleccionados-derivados-de-la-formalizaciuxf3n}

Los nombres de teoremas siguen \nolinkurl{PaperIII.PaperImprovements} y el
paquete independiente \nolinkurl{Contrib.*}. Su estado de release se rige
por el gate consolidado de build y axiomas de la Sección 13.

\subsubsection{Corolario E.5.2 --- Coeficiente lineal
forzado}\label{corolario-e.5.2-coeficiente-lineal-forzado-1}

El certificado más débil \nolinkurl{Byproduct_leading_constant_sharp},
que se sigue de inmediato de la forma normal exacta, no se cataloga por
separado.

\subsubsection{Corolario 10.2a --- Benchmark complete-split
exacto}\label{corolario-10.2a-benchmark-complete-split-exacto}

El desarrollo Lean registra la cota superior, la cota inferior y la
igualdad como \nolinkurl{Byproduct_completeSplit_cp_le_exact_value},
\nolinkurl{Byproduct_completeSplit_cp_ge_exact_value} y
\nolinkurl{Byproduct_completeSplit_cp_exact_value}. El Apéndice E.5
conserva solo la forma normal exacta y el enunciado, lógicamente
distinto, sobre el coeficiente lineal forzado.

\subsubsection{E.6.3 Interfaz uniforme de
nibble}\label{e.6.3-interfaz-uniforme-de-nibble-1}

La biblioteca independiente adjunta también contiene módulos sobre
concentración, hamiltonicidad, flujo máximo/corte mínimo, conteo por
codegrados, colas hipergeométricas y teoría de diseños. La Sección 11.7
registra su alcance sin presentarlos como corolarios adicionales de este
artículo.

\subsubsection{Corolarios en forma de
empaquetamiento}\label{corolarios-en-forma-de-empaquetamiento}

Los corolarios siguientes reformulan resultados de las Secciones 5 y 7
directamente en lenguaje de empaquetamientos. Están formalizados como
\nolinkurl{factorization_assignment_packing},
\nolinkurl{double_factorization_packing} y
\nolinkurl{reserved_gain_packing_bound_subset}. Los tres son
incondicionales: ninguno usa el Teorema 2.1 ni el Teorema 2.3.

\subsection{Referencias}\label{referencias}

{[}1{]} P. Allen, J. Böttcher, O. Cooley, and R. Mycroft, ``Tight cycles
and regular slices in dense hypergraphs,'' \emph{Journal of
Combinatorial Theory, Series A} \textbf{149} (2017), 30--100.

{[}2{]} B. Barber, D. Kühn, A. Lo, and D. Osthus, ``Edge-decompositions
of graphs with high minimum degree,'' \emph{Advances in Mathematics}
\textbf{288} (2016), 337--385.

{[}3{]} T. F. Bloom, ``Erdős Problem \#81,'' \emph{Erdős Problems},
\url{https://www.erdosproblems.com/81}, accessed July 6, 2026.

{[}4{]} O. V. Borodin, A. V. Kostochka, and D. R. Woodall, ``List edge
and list total colourings of multigraphs,'' \emph{Journal of
Combinatorial Theory, Series B} \textbf{71} (1997), no. 2, 184--204.

{[}5{]} G.-T. Chen, P. Erdős, and E. T. Ordman, ``Clique partitions of
split graphs,'' in \emph{Combinatorics, Graph Theory, Algorithms and
Applications} (Beijing, 1993), World Scientific, 1994, pp.~21--30.

{[}6{]} G. A. Dirac, ``Some theorems on abstract graphs,''
\emph{Proceedings of the London Mathematical Society} (3) \textbf{2}
(1952), 69--81.

{[}7{]} F. Dross, ``Fractional triangle decompositions in graphs with
large minimum degree,'' \emph{SIAM Journal on Discrete Mathematics}
\textbf{30} (2016), no. 1, 36--42.

{[}8{]} P. Erdős, E. T. Ordman, and Y. Zalcstein, ``Clique partitions of
chordal graphs,'' \emph{Combinatorics, Probability and Computing}
\textbf{2} (1993), no. 4, 409--415.

{[}9{]} D. Gale and L. S. Shapley, ``College admissions and the
stability of marriage,'' \emph{The American Mathematical Monthly}
\textbf{69} (1962), no. 1, 9--15.

{[}10{]} F. Galvin, ``The list chromatic index of a bipartite
multigraph,'' \emph{Journal of Combinatorial Theory, Series B}
\textbf{63} (1995), no. 1, 153--158.

{[}11{]} P. E. Haxell and V. Rödl, ``Integer and fractional packings in
dense graphs,'' \emph{Combinatorica} \textbf{21} (2001), 13--38.

{[}12{]} P. Keevash, ``The existence of designs,'' arXiv:1401.3665,
revised 2018.

{[}13{]} D. König, ``Über Graphen und ihre Anwendung auf
Determinantentheorie und Mengenlehre,'' \emph{Mathematische Annalen}
\textbf{77} (1916), 453--465.

{[}14{]} M. Krivelevich, ``On a conjecture of Tuza about packing and
covering of triangles,'' \emph{Discrete Mathematics} \textbf{142}
(1995), 281--286.

{[}15{]} J. P. Traverso Gianini, ``Reducción afín de perfiles para
empaquetamientos fraccionales de triángulos en grafos split'' (Paper I
de la serie), preprint, julio de 2026.

{[}16{]} J. P. Traverso Gianini, ``Extremizadores complete-split para un
funcional fraccional de cobertura de triángulos en grafos cordales''
(Paper II de la serie), preprint, julio de 2026.

{[}17{]} R. Yuster, ``Integer and fractional packing of families of
graphs,'' \emph{Random Structures \& Algorithms} \textbf{26} (2005),
110--118.

\end{document}
